\begin{helper}
Let \(s,k,j\in\mathbb N\) with \(1\le j\le s-1\). Then, after viewing the
natural-number exponent as a real number,
\[
  (s-1)(s-j+k)-\binom{s-j+1}{2}
  =
  \frac{s^2}{2}+sk-\frac{3s}{2}-k+\frac{3j}{2}-\frac{j^2}{2}.
\]
\end{helper}
\begin{proof}
The inequalities \(1\le j\le s-1\) imply \(1\le s\) and
\[
  s-j+1\le 2(s-1).
\]
Consequently
\[
  \binom{s-j+1}{2}
  =\frac{(s-j+1)(s-j)}2
  \le (s-1)(s-j+k).
\]
Thus the natural subtraction on the left is not truncated, so it agrees with
ordinary subtraction after casting to \(\mathbb R\).

Using
\[
  \binom{s-j+1}{2}=\frac{(s-j)(s-j+1)}2,
\]
we compute in \(\mathbb R\):
\[
\begin{aligned}
  (s-1)(s-j+k)-\binom{s-j+1}{2}
  &=(s-1)(s-j+k)-\frac{(s-j)(s-j+1)}2  \\
  &=\frac{s^2}{2}+sk-\frac{3s}{2}-k+\frac{3j}{2}-\frac{j^2}{2}.
\end{aligned}
\]
This is the claimed identity.
\end{proof}
